\documentclass[11pt,a4paper]{article}

% ------------------------------------------------------------------
%  Lecture Notes: When Does EIT Fail? A No-Go/Go Theory for
%  Arbitrary Materials  (undergraduate version)
%  Based on: "Generalized EIT No-Go/Go Theory Applicable to
%  Arbitrary Materials", v6.2 (2026-07-13)
% ------------------------------------------------------------------

\usepackage[margin=25mm]{geometry}
\usepackage{amsmath,amssymb,amsthm,mathtools}
\usepackage{bm}
\usepackage{graphicx}
\usepackage{booktabs}
\usepackage{enumitem}
\usepackage[colorlinks=true,linkcolor=blue!60!black,citecolor=blue!60!black]{hyperref}
\usepackage{tcolorbox}
\tcbuselibrary{skins,breakable}

% ---------- theorem environments ----------
\newtheorem{theorem}{Theorem}[section]
\newtheorem{lemma}[theorem]{Lemma}
\newtheorem{corollary}[theorem]{Corollary}
\newtheorem{definition}[theorem]{Definition}
\theoremstyle{remark}
\newtheorem{remark}[theorem]{Remark}
\newtheorem{example}[theorem]{Example}
\newtheorem{exercise}[theorem]{Exercise}

% ---------- boxes ----------
\newtcolorbox{keyidea}[1][Key idea]{
  breakable,enhanced,colback=blue!4,colframe=blue!50!black,
  title=#1,fonttitle=\bfseries}
\newtcolorbox{warningbox}[1][Common misconception]{
  breakable,enhanced,colback=red!4,colframe=red!55!black,
  title=#1,fonttitle=\bfseries}
\newtcolorbox{physbox}[1][Physical picture]{
  breakable,enhanced,colback=green!4,colframe=green!35!black,
  title=#1,fonttitle=\bfseries}

% ---------- macros ----------
\newcommand{\Hg}{\mathcal{H}_g}
\newcommand{\He}{\mathcal{H}_e}
\newcommand{\Hr}{\mathcal{H}_r}
\newcommand{\dd}{\mathrm{d}}
\newcommand{\im}{\mathrm{i}}
\newcommand{\Tr}{\operatorname{Tr}}
\newcommand{\rank}{\operatorname{rank}}
\newcommand{\adj}{\operatorname{adj}}
\newcommand{\spann}{\operatorname{span}}
\newcommand{\ket}[1]{|#1\rangle}
\newcommand{\bra}[1]{\langle #1|}
\newcommand{\braket}[2]{\langle #1|#2\rangle}
\newcommand{\kett}[1]{|#1\rangle\!\rangle}
\newcommand{\braa}[1]{\langle\!\langle #1|}
\newcommand{\dchi}{\delta\chi_{S}}
\newcommand{\chifull}{\chi_{\mathrm{full}}}
\newcommand{\chicut}{\chi^{(S)}_{\mathrm{cut}}}

\title{\bfseries When Does Electromagnetically Induced Transparency Fail?\\[4pt]
\Large A Rigorous No-Go/Go Theory for Arbitrary Materials\\[4pt]
\large Lecture Notes for Undergraduates}
\author{Quantum Sensing Group}
\date{July 2026}

\begin{document}
\maketitle

\begin{abstract}
Electromagnetically induced transparency (EIT) is often introduced through the
idealized three-level $\Lambda$ system, where a control laser renders an opaque
medium transparent to a probe laser. In real materials --- color centers in
diamond, alkali vapors, group-IV defects --- EIT sometimes works beautifully and
sometimes fails completely, even when a $\Lambda$-shaped level scheme exists.
These notes develop, at an undergraduate level but with full proofs, a general
theory that decides \emph{when and why} EIT is forbidden (``no-go''), suppressed,
or permitted (``go'') in an arbitrary finite-dimensional Markovian quantum
emitter. The central lesson is that a no-go statement is \emph{not}
``the susceptibility is zero''; it is the statement that a specific,
physically chosen \emph{coherence pathway} contributes nothing to the probe
response. We prove the main theorems (dark-state rank, stationary Lindblad dark
states, the Schur-complement susceptibility, the moment/Krylov exact-zero
criterion, and symmetry-protected transfer zeros) and apply them to three
concrete platforms: the nitrogen-vacancy (NV$^-$) center in diamond, alkali
atoms (Rb), and group-IV centers (SiV$^-$, SnV$^-$).
\end{abstract}

\tableofcontents

% ==================================================================
\section{Introduction: the textbook picture and why it is not enough}
% ==================================================================

\subsection{The three-level $\Lambda$ system}

Consider two long-lived lower states $\ket{g_1},\ket{g_2}$ and one excited
state $\ket{e}$. A weak \emph{probe} laser (Rabi frequency $\Omega_p$) couples
$\ket{g_1}\leftrightarrow\ket{e}$ and a strong \emph{control} laser
($\Omega_c$) couples $\ket{g_2}\leftrightarrow\ket{e}$. In the rotating frame,
with probe detuning $\Delta_p$ and two-photon detuning $\delta_2$, the standard
result for the probe optical coherence is
\begin{equation}
\rho^{(1)}_{eg_1}
 = \frac{\im\Omega_p}{2}\,
   \frac{\gamma_{21}-\im\delta_2}
        {(\gamma_{31}-\im\Delta_p)(\gamma_{21}-\im\delta_2)+|\Omega_c|^2/4}.
\label{eq:textbook}
\end{equation}
At exact two-photon resonance ($\delta_2=0$) and with an ideal ground
coherence ($\gamma_{21}\to0$), the response vanishes: the medium becomes
\emph{transparent}. The physical mechanism is destructive interference: the
population is optically pumped into the \emph{dark state}
\begin{equation}
\ket{D}\propto \Omega_c\ket{g_1}-\Omega_p\ket{g_2},
\qquad
\bra{e}H_{\mathrm{drive}}\ket{D}=0 .
\end{equation}

\subsection{Why real materials are harder}

In a real emitter there are many excited levels (fine structure, spin--orbit
splitting, phonon sidebands), many dissipation channels (radiative decay,
intersystem crossing, orbital hopping, spectral diffusion), and material-specific
selection rules. Three questions then arise, and the textbook picture cannot
answer any of them:

\begin{enumerate}[label=(Q\arabic*)]
\item Does an \emph{optically dark superposition} of the lower states even exist?
\item If it exists, is it a \emph{stationary} state of the full dissipative
      dynamics, or does dissipation destroy it?
\item Even if no exact dark state survives, can a \emph{coherent Raman pathway}
      through the excited manifold still produce a transparency feature ---
      and how strongly is it suppressed by fast dissipation?
\end{enumerate}

\begin{keyidea}
The theory in these notes classifies not \emph{materials} but
\emph{configurations}: a tuple
\[
C=\bigl(\text{material},\,\Hg,\,\He,\,
\text{lower-state pair},\,\text{polarizations},\,
\text{fields},\,T,\,\text{observable},\,S\bigr),
\]
where $S$ is the targeted long-lived coherence sector. Within one and the same
material (say NV$^-$), a spin-$\Lambda$ channel at zero field, a
transverse-field spin-$\Lambda$ channel, and an orbital-$\Lambda$ channel are
\emph{different} classification targets with different answers.
\end{keyidea}

\begin{warningbox}
A common error is to define ``EIT is forbidden'' as ``the total probe
susceptibility vanishes,'' $\chifull=0$. This is exactly backwards:
Eq.~\eqref{eq:textbook} shows that $\chifull=0$ is what \emph{perfect EIT}
looks like! We will define a no-go instead through a \emph{sector-resolved
correction} $\dchi$ --- the part of the response mediated by the targeted
ground-state coherence --- and say the channel is a no-go when
$\dchi\equiv0$ identically over a parameter domain.
\end{warningbox}

\subsection{Assumptions and scope}

Throughout, we assume:
\begin{enumerate}[label=(A\arabic*)]
\item \textbf{Finite dimensionality}: the Hilbert space (or the retained
      Liouville response space) is finite dimensional.
\item \textbf{Markovian GKSL dynamics}: the zeroth-order dynamics (control
      field, static fields, dissipation) is generated by a time-independent
      Lindblad (GKSL) generator $\mathcal{L}_c$ in a stationary rotating frame.
\item \textbf{Weak probe}: the response is computed to first order in the
      probe amplitude; the control is treated nonperturbatively.
\item \textbf{Unique steady state / well-defined group inverse} on the
      trace-zero subspace.
\end{enumerate}
Genuinely non-Markovian baths, strong-probe saturation, and
propagation-dominated collective effects are outside the main theorems and
require separate extensions.

% ==================================================================
\section{Open-system setup}
% ==================================================================

\subsection{Hilbert-space decomposition and the optical-coherence block}

Split the total Hilbert space as
\begin{equation}
\mathcal{H}=\He\oplus\Hg\oplus\Hr,
\end{equation}
where $\He$ is the operative excited manifold, $\Hg$ the lower manifold, and
$\Hr$ collects auxiliary states (singlets, other charge states, loss
channels). The probe response lives in the rectangular
\emph{optical-coherence block}
\begin{equation}
\sigma \;=\; P_e\,\rho\,P_g \;\in\; \mathcal{B}(\Hg,\He),
\end{equation}
where $P_{e},P_{g}$ are projectors onto $\He,\Hg$.

\subsection{A crucial lemma: excited-state jumps damp optical coherences}

\begin{lemma}[Excited-confined jumps act as pure damping on $\sigma$]
\label{lem:jump}
Let $L_\mu=P_e L_\mu P_e$ be a jump operator confined to the excited manifold.
Then the Lindblad dissipator
$\mathcal{D}[L_\mu](\rho)=L_\mu\rho L_\mu^\dagger
 -\tfrac12\{L_\mu^\dagger L_\mu,\rho\}$
acts on the optical-coherence block as
\begin{equation}
P_e\,\mathcal{D}[L_\mu](\rho)\,P_g
 = -\tfrac12\,L_\mu^\dagger L_\mu\,\sigma .
\end{equation}
\end{lemma}

\begin{proof}
The recycling term gives
$P_e L_\mu\rho L_\mu^\dagger P_g
 = L_\mu (P_e\rho P_e) (P_e L_\mu^\dagger P_g)=0$,
since $L_\mu^\dagger$ maps $\He\to\He$ and hence $P_e L_\mu^\dagger P_g=0$.
Of the anticommutator, only the left action survives on the $e$--$g$ block:
$P_e L_\mu^\dagger L_\mu \rho P_g=L_\mu^\dagger L_\mu\,\sigma$ and
$P_e\rho L_\mu^\dagger L_\mu P_g=0$. 
\end{proof}

\begin{physbox}
Even a jump that merely \emph{shuffles population} among excited branches
(e.g.\ phonon-induced orbital hopping $\ket{E_x}\leftrightarrow\ket{E_y}$ in
NV$^-$) acts on every optical coherence as \emph{loss}. There is no
``recycling'' in the rectangular block. This is why fast phonon dynamics can
kill EIT even without removing any population from the excited manifold ---
and it is the microscopic origin of the fast-damping expansions in
Section~\ref{sec:moments}.
\end{physbox}

\begin{example}[Rate mapping --- a frequent factor-of-4 trap]
\label{ex:rates}
For symmetric hopping $\ket{X}\leftrightarrow\ket{Y}$ at rate $k$ each way,
the population imbalance obeys
$\dot{(p_X-p_Y)}=-\Gamma_{XY}(p_X-p_Y)$ with $\Gamma_{XY}=2k$.
But by Lemma~\ref{lem:jump}, each optical coherence $\sigma_{Xg}$ acquires
damping $k/2=\Gamma_{XY}/4$. If a paper quotes ``the orbital relaxation rate
$\Gamma_{XY}$,'' the number entering the optical-coherence equation is
$\Gamma_{XY}/4$, not $\Gamma_{XY}$. Mixing these conventions changes every
quantitative prediction by factors of 2--4.
\end{example}

\subsection{The generator of the optical coherence}

Collecting the Hamiltonian commutator and Lemma~\ref{lem:jump}, the undriven
optical coherence obeys
\begin{equation}
\dot\sigma
 = -\im\bigl(H_e\sigma-\sigma H_g\bigr)
   -\tfrac12 R_e\,\sigma + \text{(corrections)} + \text{sources},
\qquad
R_e=\sum_\mu L_\mu^\dagger L_\mu .
\end{equation}
Vectorizing with $\operatorname{vec}(AXB)=(B^{\mathsf T}\!\otimes A)
\operatorname{vec}X$ gives a matrix generator $A(z)$ acting on
$\operatorname{vec}\sigma$. When the lower index is not mixed (no ground-state
relaxation into the block, no correlated jumps), the problem reduces to an
excited-manifold resolvent equation
\begin{equation}
A_j(z)\,\ket{x_j}=\ket{s_j},
\qquad G(z)=A(z)^{-1}.
\end{equation}

% ==================================================================
\section{Dark states: two theorems that must not be confused}
% ==================================================================

\subsection{The optical dark subspace}

In the rotating frame, the drive couples lower to excited states through
\begin{equation}
V_+ = P_e H_{\mathrm{drive}} P_g=\tfrac12\Omega,
\qquad
\Omega:\Hg\to\He,
\qquad
\Omega=(\Omega_1 d_1,\dots,\Omega_{N_g} d_{N_g}),
\quad (d_a)_j=\bra{e_j}\bm{\epsilon}_a\!\cdot\!\hat{\bm d}\ket{g_a},
\end{equation}
where the \emph{complex} phases of the dipole vectors $d_a$ are retained.

\begin{theorem}[1A: optical dark-subspace rank]
\label{thm:1A}
A lower-state vector $\ket{v}\in\Hg$ is instantaneously decoupled from optical
excitation, $V_+\ket{v}=0$, if and only if $\Omega\ket{v}=0$. Hence the
optical dark subspace is $\mathcal{D}_{\mathrm{opt}}=\ker\Omega$ with
\begin{equation}
\dim\mathcal{D}_{\mathrm{opt}} = N_g-\rank\Omega .
\end{equation}
\end{theorem}

\begin{proof}
$V_+=\Omega/2$ gives the equivalence; the dimension formula is the
rank--nullity theorem. Note $\Omega\mapsto U_e\Omega U_g^\dagger$ under basis
changes, so rank and nullity are basis invariant.
\end{proof}

\begin{corollary}
For two lower states with both Rabi frequencies nonzero, a nontrivial dark
vector exists iff $\rank(d_1,d_2)<2$, i.e.\ (for nonzero dipoles) iff the two
dipole vectors are \emph{parallel} in $\He$: $d_1\parallel d_2$.
\end{corollary}

\begin{physbox}
This is the first, purely geometric filter: EIT interference requires the two
optical legs to reach the \emph{same} superposition of excited states. In an
alkali atom driving both legs to a single hyperfine excited level, $d_1$ and
$d_2$ are trivially parallel (both are scalars). In NV$^-$ or SiV$^-$, the two
legs may reach \emph{different} excited orbitals, and the rank condition
already becomes nontrivial.
\end{physbox}

\subsection{Stationary Lindblad dark states}

An optical dark vector need not survive dissipation. The full condition is:

\begin{theorem}[1B: stationary pure dark state of a Lindblad equation]
\label{thm:1B}
Let $\dot\rho=\mathcal{L}(\rho)=-\im[H,\rho]
+\sum_\mu\bigl(L_\mu\rho L_\mu^\dagger-\tfrac12\{L_\mu^\dagger L_\mu,\rho\}\bigr)$.
For a normalized $\ket{D}$, the pure state $\rho_D=\ket{D}\bra{D}$ is
stationary iff there exist $\lambda_\mu\in\mathbb C$ and $h\in\mathbb R$ with
\begin{equation}
L_\mu\ket{D}=\lambda_\mu\ket{D}\ \ (\forall\mu),
\qquad
\Bigl[H+\tfrac{\im}{2}\sum_\mu(\lambda_\mu^*L_\mu-\lambda_\mu L_\mu^\dagger)\Bigr]
\ket{D}=h\ket{D}.
\end{equation}
\end{theorem}

\begin{proof}[Proof sketch]
For a stationary pure state, the purity is constant. A direct computation
gives
\[
\frac{\dd}{\dd t}\Tr(\rho_D^2)
 = -2\sum_\mu\Bigl(\langle L_\mu^\dagger L_\mu\rangle_D
   -|\langle L_\mu\rangle_D|^2\Bigr),
\]
a sum of nonnegative variances. Stationarity forces each variance to vanish,
i.e.\ $\ket{D}$ is an eigenvector of \emph{every} jump operator:
$L_\mu\ket{D}=\lambda_\mu\ket{D}$. Substituting back and projecting with
$Q_D=I-\ket{D}\bra{D}$ yields
$Q_D\bigl[H+\tfrac{\im}{2}\sum_\mu(\lambda_\mu^*L_\mu
-\lambda_\mu L_\mu^\dagger)\bigr]\ket{D}=0$; the bracketed operator is
Hermitian, so its action on $\ket{D}$ is parallel to $\ket{D}$ with a real
eigenvalue $h$. Conversely, inserting both conditions into $\mathcal{L}(\rho_D)$
kills all off-diagonal terms and, using trace preservation, gives
$\mathcal{L}(\rho_D)=0$.
\end{proof}

For an \emph{EIT} dark state one additionally imposes $P_e\ket{D}=0$ and
$\Omega\ket{D}=0$. Spontaneous emission from excited states satisfies
$L_\mu\ket{D}=0$ automatically; but ground-state dephasing, ground-state
relaxation, or microwave noise generically \emph{violate} the
common-eigenvector condition. Then:

\begin{warningbox}
An optical dark vector can exist ($\ker\Omega\neq\{0\}$, Theorem~\ref{thm:1A})
while \emph{no stationary dark state exists} (Theorem~\ref{thm:1B} fails). The
two theorems answer different questions. Only the pair together, plus the
response analysis below, decides observable EIT.
\end{warningbox}

\subsection{Fixed terminology}

\begin{center}
\begin{tabular}{ll}
\toprule
\textbf{Notion} & \textbf{Defining condition}\\
\midrule
optical dark capable & $\dim\ker\Omega>0$\\
Hamiltonian dark eigenstate & $\Omega\ket{D}=0$ and $H_g^{\mathrm{(rot)}}\ket{D}=E_D\ket{D}$\\
stationary Lindblad dark state & Theorem~\ref{thm:1B}\\
observable EIT & narrow coherent transparency in the full response\\
\bottomrule
\end{tabular}
\end{center}

% ==================================================================
\section{The coherent-response kernel and the exact susceptibility}
% ==================================================================

\subsection{The kernel $K(z)=C^\dagger G(z)\,C$}

With the dipole matrix $C=(d_1,\dots,d_{N_g})$ and the Green operator
$G(z)=A(z)^{-1}$ of the optical-coherence generator, define
\begin{equation}
K_{ab}(z)=d_a^\dagger\,G(z)\,d_b .
\end{equation}
The diagonal $S_a=K_{aa}$ is the single-leg optical self-energy; the
off-diagonal $K_{ab}$ is the \emph{coherent Raman vertex} mediated by the
excited manifold. These are invariant under unitary basis changes in $\He$.
For a non-normal generator, note $K_{21}\neq K_{12}^*$ in general --- both
must be computed.

\subsection{The exact weak-probe susceptibility}

To first order in the probe, the full-Liouvillian response is
\begin{equation}
\rho^{(1)}=-\mathcal{L}_c^{\#}\,\mathcal{V}_p\,\rho^{(0)},
\qquad
\chi_p=-\chi_0\,\braa{O_p}\mathcal{L}_c^{\#}\kett{\mathcal{V}_p\rho^{(0)}},
\label{eq:fullresp}
\end{equation}
where $\mathcal{L}_c^{\#}$ is the group inverse on the trace-zero subspace
(the Liouvillian has a zero mode --- the steady state --- so an ordinary
inverse does not exist; this is not a technicality, it is a definition).

For a closed two-lower-state block the first-order variables are the optical
coherence vector $x\in\He$ and the ground coherence
$s=\rho^{(1)}_{g_2g_1}$, obeying
\begin{equation}
\begin{pmatrix}
A & \im\Omega_c d_2/2\\
\im\Omega_c^* d_2^\dagger/2 & \gamma_g
\end{pmatrix}
\begin{pmatrix} x\\ s\end{pmatrix}
=
\begin{pmatrix} \im\Omega_p d_1/2\\ 0\end{pmatrix},
\qquad
\gamma_g=\gamma_{12}-\im\delta_2 .
\label{eq:block}
\end{equation}

\begin{theorem}[2A: reduced susceptibility as an exact Schur complement]
\label{thm:2A}
Let $G=A^{-1}$, $S_i=d_i^\dagger G d_i$, $K_{ij}=d_i^\dagger G d_j$,
$\beta=|\Omega_c|^2/4$, and define the normalized probe response
$\Xi=d_1^\dagger x/(\im\Omega_p/2)$. Then, whenever
$\gamma_g+\beta S_2\neq0$,
\begin{equation}
\boxed{\;\Xi \;=\; S_1-\frac{\beta\,K_{12}K_{21}}{\gamma_g+\beta S_2}\;}
\label{eq:Xi}
\end{equation}
with no adiabatic approximation.
\end{theorem}

\begin{proof}
The first row of \eqref{eq:block} gives
$x=G\,\tfrac{\im\Omega_p}{2}d_1-G\,\tfrac{\im\Omega_c}{2}d_2\,s$.
Insert into the second row \emph{without dividing by $\gamma_g$}:
\[
(\gamma_g+\beta S_2)\,s
 = -\frac{\im\Omega_c^*}{2}\,d_2^\dagger G\,\frac{\im\Omega_p}{2}\,d_1
 = \frac{\Omega_c^*\Omega_p}{4}\,K_{21},
\]
so $s=\tfrac{\Omega_c^*\Omega_p}{4}K_{21}/(\gamma_g+\beta S_2)$. Substituting
back into $x$, projecting with $d_1^\dagger$, and dividing by $\im\Omega_p/2$
gives \eqref{eq:Xi}. This is precisely the Schur complement of the block matrix
in \eqref{eq:block}; because we never divided by $\gamma_g$ itself, the ideal
limit $\gamma_g\to0$ is included.
\end{proof}

\begin{example}[Sanity check: the textbook $\Lambda$ system]
\label{ex:textbook}
One excited state, scalar dipoles $d_1=d_2=1$, $A=\gamma_{31}-\im\Delta_p$,
$\gamma_g=\gamma_{21}-\im\delta_2$. Then $S_1=S_2=K_{12}=K_{21}=A^{-1}$ and
\[
\Xi=\frac1A-\frac{\beta/A^2}{\gamma_g+\beta/A}
   =\frac{\gamma_g}{A\gamma_g+\beta},
\]
which is exactly Eq.~\eqref{eq:textbook}. The general machinery contains the
textbook result as a one-line special case.
\end{example}

\begin{example}[Multi-excited-state (Mishina-type) media]
For $N_e$ resolved excited branches with diagonal generator
$A=\operatorname{diag}(a_j)$, $a_j=\gamma_j-\im\Delta_j$,
\[
S_i=\sum_j\frac{|d_{ij}|^2}{a_j},
\qquad
K_{12}=\sum_j\frac{d_{1j}^*d_{2j}}{a_j},
\]
so the Raman vertex is a \emph{coherent sum of amplitudes} through all
branches --- pathways can interfere destructively even when every individual
dipole is nonzero. Inhomogeneous (Doppler, strain) broadening is applied
\emph{after} computing the local response:
$\bar\chi=\int\dd\lambda\,P(\lambda)\chi_\lambda$. Averaging the Hamiltonian
first is generally wrong.
\end{example}

% ==================================================================
\section{Defining the no-go correctly: the sector-resolved correction}
% ==================================================================

\subsection{Cutting the pathway, not the control}

Fix a physically specified \emph{long-lived sector} $S$ (for us: the
ground-state coherence $\rho_{g_2g_1}$) and its complement $X$ (the direct
optical response). Writing the first-order block as
\begin{equation}
M=\begin{pmatrix} A & B\\ C & G_g\end{pmatrix},
\qquad
\kett{b_p}=\begin{pmatrix}b_p\\0\end{pmatrix},
\qquad
\braa{d_p}=(d^\dagger\ \ 0),
\end{equation}
define
\begin{equation}
\chifull=\chi_0\,d^\dagger x,
\qquad
\chicut=\chi_0\,d^\dagger A^{-1}b_p,
\qquad
\boxed{\;\dchi\equiv\chifull-\chicut\;}
\end{equation}
where the \emph{cut} sets $B=C=0$ while keeping $A$, $G_g$, source, detector
and the zeroth-order state fixed. This is an \emph{algebraic counterfactual}:
it is \emph{not} the control-off spectrum (switching the control off would
also change optical pumping, light shifts, and the steady state).

\begin{theorem}[2B: sector-resolved EIT no-go]
\label{thm:2B}
If $A$ and $S_g=G_g-CA^{-1}B$ are invertible at a regular point, then
\begin{equation}
\dchi=\chi_0\,d^\dagger A^{-1} B\, S_g^{-1} C A^{-1} b_p .
\end{equation}
If on a connected analytic parameter domain $\dchi(\theta)\equiv0$, the net
$S$-mediated contribution to the response vanishes identically and the
configuration is an \emph{exact $S$-sector EIT no-go}. For the scalar block
\eqref{eq:block} with nonzero control and finite denominator,
\begin{equation}
\dchi=0\iff K_{12}K_{21}=0 .
\end{equation}
\end{theorem}

\begin{proof}
From \eqref{eq:block}--\eqref{eq:Xi},
$\delta\Xi_S=\Xi_{\mathrm{full}}-\Xi_{\mathrm{cut}}
=-\beta K_{12}K_{21}/(\gamma_g+\beta S_2)$;
with $\beta\neq0$ and a finite denominator this vanishes iff
$K_{12}K_{21}=0$. The matrix form follows from the block inverse
$x=A^{-1}b_p+A^{-1}BS_g^{-1}CA^{-1}b_p$ (derived exactly as in
Theorem~\ref{thm:2A}, valid even for singular $G_g$).
\end{proof}

\subsection{Why $\chifull=0$ must never be the criterion}

\begin{example}[The ideal $\Lambda$ counterexample]
Take $A=1$, $d_1=d_2=1$, $\beta=1$, $\gamma_g=0$. Then
$S_1=S_2=K_{12}=K_{21}=1$ and
\[
\Xi_{\mathrm{cut}}=1,\qquad
\Xi_{\mathrm{full}}=1-\frac{1}{1}=0,\qquad
\delta\Xi_S=-1\neq0 .
\]
The total response is zero --- \emph{perfect transparency}, ideal EIT --- while
the sector correction is maximally nonzero. A zero of $\chifull$ is a
\emph{go} signature, not a no-go.
\end{example}

\begin{warningbox}
Three logically distinct kinds of ``zero'' must be reported separately:
\begin{enumerate}
\item \textbf{Symmetry-protected zero}: enforced by reducing projectors of the
      full generator; robust against symmetry-preserving perturbations.
\item \textbf{Krylov exact zero}: every finite transfer path is orthogonal to
      the readout (Theorem~\ref{thm:krylov}); can occur without any explicit
      symmetry.
\item \textbf{Accidental parameter zero}: cancellation at one detuning or
      field value; \emph{not} a domain-wide no-go.
\end{enumerate}
A numerical zero at finitely many frequencies proves nothing about an
identically vanishing transfer.
\end{warningbox}

% ==================================================================
\section{Fast damping: the moment hierarchy and exact Krylov zeros}
\label{sec:moments}
% ==================================================================

\subsection{Setup}

Decompose the optical-coherence generator as
\begin{equation}
A_\Gamma(z)=\Gamma D + A_0(z),
\end{equation}
where $\Gamma>0$ is the fast-rate scale (e.g.\ phonon-induced orbital hopping
at elevated temperature), $D$ the dimensionless fast-generator shape, and
$A_0$ collects detunings, fine structure, slow losses and fields.
(Remember Example~\ref{ex:rates}: $\Gamma$ must be read off from the
optical-coherence equation, not from a population-relaxation rate.)

Define, assuming $D$ invertible,
\begin{equation}
X(z)=-D^{-1}A_0(z),
\qquad
\nu=D^{-1}c,
\qquad
K_\Gamma(z)=p^\dagger A_\Gamma^{-1}c
          =p^\dagger(\Gamma I-X)^{-1}\nu .
\end{equation}

\subsection{The moment expansion}

For $q=\|X\|/\Gamma<1$ the Neumann series gives
\begin{equation}
K_\Gamma=\sum_{n=0}^\infty\frac{M_n}{\Gamma^{n+1}},
\qquad
M_n=p^\dagger X^n\nu .
\end{equation}

\begin{theorem}[3: first nonzero moment fixes the suppression order]
If $M_0=\dots=M_{m-1}=0$ and $M_m\neq0$, then as $\Gamma\to\infty$
\begin{equation}
K_\Gamma=\frac{M_m}{\Gamma^{m+1}}+O(\Gamma^{-m-2}).
\end{equation}
\end{theorem}

\begin{physbox}
Interpretation as a \emph{sector graph}: split the coherence space into
sectors preserved by $D$ and the diagonal part of $A_0$, and treat the
off-diagonal coupling $V$ as edges. If the minimal number of $V$-insertions
connecting the source sector to the readout sector is $d$, then $M_n=0$ for
$n<d$ and $K_\Gamma=O(\Gamma^{-1-d})$. Each ``hop'' across a sector boundary
costs one power of $1/\Gamma$. Fast dissipation acts as a
\emph{dissipative selection rule}.
\end{physbox}

\subsection{Exact zeros: the Krylov/Cayley--Hamilton theorem}

\begin{theorem}[4: exact finite-dimensional transfer zero]
\label{thm:krylov}
Let the response space have dimension $N$ and let $D$ be invertible at fixed
$z$. The following are equivalent:
\begin{enumerate}[label=(\roman*)]
\item $M_n=p^\dagger X^n\nu=0$ for $n=0,\dots,N-1$;
\item $p\perp\mathcal K_N(X,\nu)
      =\spann\{\nu,X\nu,\dots,X^{N-1}\nu\}$;
\item $p^\dagger q(X)\nu=0$ for every polynomial $q$;
\item $K_\Gamma=0$ for every $\Gamma\notin\operatorname{spec}X$.
\end{enumerate}
\end{theorem}

\begin{proof}
(i)$\Leftrightarrow$(ii) is the definition of the Krylov space. By
Cayley--Hamilton, $X^N$ is a polynomial in $I,X,\dots,X^{N-1}$; hence (i)
implies inductively that \emph{all} moments vanish, giving (iii). For
(iii)$\Rightarrow$(iv), use the adjugate representation
\[
(\Gamma I-X)^{-1}=\frac{\adj(\Gamma I-X)}{\det(\Gamma I-X)},
\]
where the adjugate is a polynomial in $X$ of degree $\le N-1$; by (iii) its
matrix element between $p$ and $\nu$ vanishes. Conversely, if (iv) holds on an
open set of large $\Gamma$, every Laurent coefficient in $1/\Gamma$ vanishes,
which is (i).
\end{proof}

\begin{keyidea}
Only \emph{finitely many} moments ($N$ of them) need to be checked to certify
an \emph{all-orders exact zero}, and --- via the adjugate --- the conclusion
is \emph{not} restricted to the convergence domain of the $1/\Gamma$ series.
This turns ``EIT is exactly forbidden'' into a finite, machine-checkable
computation.
\end{keyidea}

\begin{remark}[What $D>0$ is actually needed for]
Theorem~\ref{thm:krylov} needs only $D$ \emph{invertible}. Positivity
($\operatorname{Re}D\ge d_{\min}I>0$) is a separate, physical condition
ensuring that large $\Gamma$ uniformly damps all modes. A non-accretive $D$
can produce large finite-$\Gamma$ responses via pseudospectral amplification
even when the asymptotic scaling holds --- check resolvent norms before
claiming visibility bounds.
\end{remark}

\subsection{Singular damping: protected channels}

If $D=D^\dagger\ge0$ is \emph{singular}, let $P=\operatorname{Proj}(\ker D)$
and $Q_D=I-P$. The Schur complement on the slow subspace gives
\begin{equation}
K_\Gamma=p^\dagger P\,(PA_0P)^{-1}P c+O(\Gamma^{-1}) :
\end{equation}
an $O(1)$ channel survives arbitrary strong damping iff \emph{both} probe and
control vectors have components in the genuine kernel of the damping operator
(and the leading scalar does not cancel). Such \emph{protected channels} are
one route to robust EIT in noisy solids.

\subsection{Symmetry-protected transfer zeros}

\begin{theorem}[5: reducing-symmetry transfer zero]
\label{thm:symm}
Suppose the response space admits orthogonal projectors $\{P_\alpha\}$,
$\sum_\alpha P_\alpha=I$, with $[A(z),P_\alpha]=0$ for all $\alpha,z$. If the
source and readout lie in different sectors, $P_\beta c=c$ and
$p^\dagger P_\alpha=p^\dagger$ with $\alpha\neq\beta$, then
$p^\dagger A(z)^{-1}c=0$ wherever $A$ is invertible.
\end{theorem}

\begin{proof}
$[A,P_\alpha]=0\Rightarrow[A^{-1},P_\alpha]=0$, so
$p^\dagger A^{-1}c
=p^\dagger P_\alpha A^{-1}P_\beta c
=p^\dagger A^{-1}P_\alpha P_\beta c=0$.
\end{proof}

The same statement lifts to the full Liouvillian with the group inverse
$\mathcal{L}_c^{\#}$, provided the sector projectors commute with
$\mathcal{L}_c$ \emph{and} with the steady-state projector. A sufficient
microscopic condition: a unitary $U$ with $UH_cU^\dagger=H_c$ and
$UL_\mu U^\dagger=\sum_\nu V_{\mu\nu}L_\nu$ ($V$ unitary, rates invariant).

\begin{warningbox}[The symmetry audit]
Before declaring a symmetry-protected no-go, one must audit \emph{every}
block: lower Hamiltonian, excited Hamiltonian including spin--orbit,
spin--spin, Zeeman and strain terms, the control Hamiltonian for the actual
polarization, all jump operators (radiative decay, ISC, orbital hopping,
ground dephasing), the steady state, the source, and the detector. If any one
item breaks the sector, the exact zero fails, and the symmetry-breaking term
must be reclassified as a perturbation $V$ feeding the moment hierarchy.
\end{warningbox}

% ==================================================================
\section{Classification taxonomy}
% ==================================================================

\begin{center}
\renewcommand{\arraystretch}{1.25}
\begin{tabular}{clp{8.6cm}}
\toprule
Class & Name & Defining condition\\
\midrule
0 & unspecified & inputs insufficient for classification\\
I & optical dark capable & $\dim\ker\Omega>0$ (Thm.~\ref{thm:1A} only)\\
II & \textbf{exact sector no-go} & $\dchi(\theta)\equiv0$ on the domain
     (proved by symmetry, Krylov, or adjugate identity);
     scalar case $\Leftrightarrow K_{12}K_{21}\equiv0$\\
III & asymptotic suppression & first nonzero moments $m_{12},m_{21}$:
     $K_{ij}=O(\Gamma^{-m_{ij}-1})$; e.g.\
     $\delta\Xi_S=O(\Gamma^{-m_{12}-m_{21}-2})$ or $O(\Gamma^{-m_{12}-m_{21}-1})$
     depending on $|\beta S_2/\gamma_g|$\\
IV & leading-order kernel go & $K_{12},K_{21}=O(\Gamma^{-1})$ with nonzero
     leading product\\
V & protected go & $O(1)$ channel via $\ker D$ Schur complement\\
VI & practical no-go & kernel nonzero but contrast below threshold
     (SNR, power, ionization, optical depth)\\
VII & observable go & narrow coherent transparency verified as EIT
     (not ATS)\\
VIII & tunable transition & fields/strain/polarization move the
     configuration among II--VII\\
\bottomrule
\end{tabular}
\end{center}

\begin{warningbox}[EIT vs Autler--Townes]
A transparency dip alone identifies nothing. A strong control splits the
excited line into an Autler--Townes doublet even with zero ground coherence.
Discriminators: dependence on $\gamma_g$, control-power scaling of the dip
width, pole/residue structure, two-photon linewidth, and the difference
between $\chifull$ and $\chicut$.
\end{warningbox}

% ==================================================================
\section{Worked platform examples}
% ==================================================================

\subsection{Alkali atoms (Rb): the canonical go}

\paragraph{Configuration.}
$^{87}$Rb D$_1$ line; lower states $\ket{g_1},\ket{g_2}$ = two hyperfine
ground states ($F=1,2$, $5{}^2S_{1/2}$); operative excited manifold = one (or
a few) hyperfine levels of $5{}^2P_{1/2}$; sector $S$ = the hyperfine
ground coherence.

\paragraph{Analysis.}
\begin{itemize}
\item \textbf{Thm.~\ref{thm:1A}}: with a single dominant excited level, both
      dipole ``vectors'' are scalars, $\rank(d_1,d_2)=1<2$: an optical dark
      state exists for any polarization pair that drives both legs.
\item \textbf{Thm.~\ref{thm:1B}}: radiative decay starts in $\He$, so
      $L_\mu\ket{D}=0$; ground-state decoherence (wall collisions, magnetic
      noise) is slow, $\gamma_{21}/2\pi\sim\mathrm{kHz}$ or below in
      buffer-gas/coated cells. The stationary-dark-state conditions are met to
      excellent approximation.
\item \textbf{Kernel}: $K_{12}=K_{21}=A^{-1}=O(\Gamma_{\mathrm{rad}}^{-1})$
      --- Class IV at leading order, and in practice Class VII: this is the
      textbook Example~\ref{ex:textbook}.
\item \textbf{Real-world corrections}: Doppler averaging must be done at the
      response level, $\bar\chi=\int\dd v\,P(v)\chi(v)$; with co-propagating
      beams the two-photon detuning is nearly Doppler-free, which is why the
      EIT window survives in hot vapor. Multiple excited hyperfine levels turn
      $K_{12}$ into a coherent sum $\sum_j d_{1j}^*d_{2j}/a_j$ whose partial
      cancellations modify (but do not kill) the contrast.
\end{itemize}

\paragraph{Verdict.} Class VII (observable go). Rb is the benchmark any
correct theory must recover --- and Theorem~\ref{thm:2A} does, exactly.

\subsection{NV$^-$ in diamond: the instructive no-go}

\paragraph{Configuration 1: zero-field ZPL spin-$\Lambda$.}
Lower states = two spin sublevels ($m_s=\pm1$ or $0,\pm1$ pairs) of the
$^3A_2$ ground triplet; excited manifold = $^3E$; probe/control on the
zero-phonon line; sector $S$ = the ground \emph{spin} coherence.

\paragraph{Analysis.}
\begin{itemize}
\item \textbf{Spin-scalar dipole.} To leading order the electric dipole is
      orbital: $\hat{\bm d}=\hat{\bm d}_{\mathrm{orb}}\otimes I_{\mathrm{spin}}$.
      For a lower pair sharing the same orbital and differing by
      \emph{orthogonal spin} states, the leading moment is a projected dipole
      overlap
      \[
      M_0\propto p^\dagger c
      =\bra{g_1}\hat d_p^\dagger\,\Pi_E\,\hat d_c\ket{g_2}
      \;\propto\;\braket{s_1}{s_2}\;=\;0 .
      \]
      The Raman vertex vanishes at leading order: $M_0=0$.
\item \textbf{Is this an exact no-go (Class II)?} Only in a simplified
      sector-preserving model. In the physical six-level $^3E$ manifold,
      transverse spin--orbit coupling, nonsecular spin--spin terms, and
      hyperfine coupling generically break the spin sector. In fact, in the
      reduced 6-level analysis the correct no-go criterion sharpens from the
      naive overlap condition $\braket{s_1}{s_2}=0$ to a \emph{commutator}
      condition $[\bar V_{\mathrm{spin}},H_g^{\mathrm{spin}}]=0$: the sector
      must be preserved by the \emph{averaged} coupling, not merely orthogonal
      at one instant.
\item \textbf{Dissipative selection rule (Class III).} What survives
      rigorously is the asymptotic statement: with phonon-induced orbital
      hopping providing full-rank fast damping $\Gamma(T)$ on the
      optical-coherence block (Lemma~\ref{lem:jump}, with the
      $\Gamma_{XY}/4$ mapping of Example~\ref{ex:rates}),
      \[
      M_0=0,\qquad K=O(\Gamma^{-2}),\qquad
      \delta\Xi_S=O(\Gamma^{-3})\ \text{or}\ O(\Gamma^{-4}),
      \]
      the branch depending on $|\beta S_2/\gamma_g|$. This is why NV spin-EIT
      degrades so violently with temperature: raising $T$ increases
      $\Gamma(T)$ and the sector-mediated correction collapses with a high
      power of $1/\Gamma$ --- a \emph{thermal collapse law} that is directly
      falsifiable by measuring the contrast against
      $x(T)=\Gamma_{\mathrm{fast}}(T)/\Delta_{\mathrm{sel}}$.
\end{itemize}

\paragraph{Configuration 2: transverse-field escape channel (Class VIII).}
Apply a small transverse magnetic field $B_\perp$. This is a sector-changing
perturbation $V_{\mathrm{mix}}=\lambda V_1$ with $\lambda\propto B_\perp$. One
insertion opens the source--readout path, so both kernels open at first order:
\[
K_{12},K_{21}=O(\lambda\Gamma^{-2}),
\qquad
K_{12}K_{21}=O(\lambda^2\Gamma^{-4}),
\qquad
C_{\mathrm{abs}}=O(B_\perp^2\Gamma^{-3})\ \text{or}\ O(B_\perp^2\Gamma^{-2}).
\]
The predicted signature $|K|^2\propto B_\perp^2$ is a clean experimental null
test: at $B_\perp=0$ the channel is (asymptotically) closed; it opens
\emph{quadratically} in the contrast as $B_\perp$ is turned on. (In the
simplified model, the even channel has $M_1\neq0$ while the odd channel has
$M_1\simeq0$; this is model-specific and must be re-audited with the full
microscopic Hamiltonian.)

\paragraph{Verdict.} Zero-field spin-$\Lambda$: Class III (asymptotic
dissipative no-go), \emph{not} an unconditional exact no-go. Transverse-field
channel: Class VIII, tunable escape with $B_\perp^2$ scaling.

\subsection{Group-IV centers (SiV$^-$, SnV$^-$): the orbital go}

\paragraph{Configuration.}
Inversion-symmetric defects with ground manifold $^2E_g$ split by spin--orbit
coupling into orbital branches (SiV: $\sim$46\,GHz; SnV: $\sim$850\,GHz), and
excited manifold $^2E_u$. Two natural channels exist \emph{in the same
material}:
\begin{enumerate}
\item \textbf{Same-spin orbital-$\Lambda$}: lower pair = two orbital branches
      with the \emph{same} spin.
\item \textbf{Spin-$\Lambda$}: lower pair = two spin states.
\end{enumerate}

\paragraph{Analysis of the orbital-$\Lambda$.}
The dipole is orbital, so it acts \emph{within} the fixed spin sector: no
spin-scalar obstruction. The leading criterion is
\[
d_1^\dagger D^{-1} d_2\neq0
\;\Longrightarrow\;
K_{12}=O(\Gamma^{-1}) :
\]
Class IV, a leading-order kernel go. Both legs connect to the same excited
orbital doublet with symmetry-allowed, generally non-parallel polarization
selection rules ($D_{3d}$ point group: check that
$\Gamma(g_1)^*\otimes\Gamma(\hat O_{pc})\otimes\Gamma(g_2)$ contains the
totally symmetric representation, with
$\hat O_{pc}=\hat d_p^\dagger\Pi_E\hat d_c$ selected by
$\bm\epsilon_p^*\otimes\bm\epsilon_c$).

\paragraph{But kernel go $\neq$ observable go.}
The orbital ground coherence is damped by single-phonon transitions between
the branches: $\gamma_g$ is large at 4\,K in SiV (orbital $T_1\sim$~tens of ns)
and the denominator $\gamma_g+\beta S_2$ eats the correction --- Class VI
(practical no-go) unless one cools to millikelvin, strains the crystal to
increase the branch splitting (suppressing the phonon density of states), or
moves to SnV, whose larger spin--orbit splitting pushes the phonon process to
higher energy and freezes it out at moderate temperatures. This is a textbook
Class VIII story: temperature and strain move the same configuration along
IV $\to$ VI $\to$ VII.

\begin{warningbox}
Orbital degeneracy alone does \emph{not} imply a go. One needs, cumulatively:
symmetry-allowed transitions, the right polarization tensor, nonzero reduced
matrix elements, nonzero spin overlap, correct manifold selection, \emph{and}
damping/coherence budgets that permit practical visibility.
\end{warningbox}

\subsection{Comparison table}

\begin{center}
\renewcommand{\arraystretch}{1.3}
\begin{tabular}{lcccc}
\toprule
Configuration & dark rank & first moment & class & key knob\\
\midrule
Rb hyperfine $\Lambda$ & $\ge1$ & $M_0\neq0$ & VII & Doppler-free geometry\\
NV$^-$ spin-$\Lambda$, $B=0$ & $\ge1$ & $M_0=0$ & III & temperature ($\Gamma(T)$)\\
NV$^-$ spin-$\Lambda$, $B_\perp\!\neq\!0$ & $\ge1$ & $M_1\propto B_\perp$ & VIII & $|K|^2\propto B_\perp^2$\\
SiV$^-$/SnV$^-$ orbital-$\Lambda$ & $\ge1$ & $M_0\neq0$ & IV/VI$\to$VII & $T$, strain, splitting\\
\bottomrule
\end{tabular}
\end{center}

% ==================================================================
\section{The classification algorithm (summary)}
% ==================================================================

\begin{enumerate}
\item Fix the configuration tuple $C$, including the sector $S$.
\item Build $H_g$, $H_e$, complex dipole vectors, all jump operators with
      correct rate conventions (Example~\ref{ex:rates}).
\item Solve $\mathcal L_c\rho^{(0)}=0$; restrict the first-order problem to
      the trace-zero space.
\item Compute $\dim\ker\Omega$ (Thm.~\ref{thm:1A}) and test
      Thm.~\ref{thm:1B} for stationarity.
\item Where valid, reduce to $A(z)$ and compute $S_1,S_2,K_{12},K_{21}$;
      otherwise keep the Liouville-space kernel.
\item Construct $\chifull$, $\chicut$, $\dchi$ (Thm.~\ref{thm:2B}).
\item Test an exact zero of the \emph{targeted} cross-transfer, in order:
      reducing symmetry (Thm.~\ref{thm:symm}, with the full audit),
      finite Krylov moments (Thm.~\ref{thm:krylov}),
      adjugate/analytic identity; at singular points use Laurent/Drazin
      analysis --- never read $0/0$ as $0$.
\item If not exact, decompose $A=\Gamma D+A_0$; if $D$ full rank, get the
      suppression order from the first nonzero moment; if singular, compute
      the protected-channel Schur complement.
\item Ensemble-average static disorder \emph{at the response level};
      convert to the measured observable; separate EIT from ATS;
      compare contrast with the experimental threshold.
\item Report: class, moment order, $K_{12},K_{21}$, $\dchi$, contrast,
      margins, and \emph{all} assumptions and projectors used.
\end{enumerate}

% ==================================================================
\section{Exercises}
% ==================================================================

\begin{exercise}
Verify by direct substitution that Example~\ref{ex:textbook} reproduces
Eq.~\eqref{eq:textbook}, and show that in the limit $\gamma_{21}=\delta_2=0$
one has $\Xi_{\mathrm{full}}=0$ but $\delta\Xi_S=-A^{-1}\neq0$.
\end{exercise}

\begin{exercise}
Prove Lemma~\ref{lem:jump} for the asymmetric case
$k_{Y\leftarrow X}\neq k_{X\leftarrow Y}$ and show that the resulting damping
matrix on the two optical coherences is
$\tfrac12\operatorname{diag}(k_{Y\leftarrow X},k_{X\leftarrow Y})$ ---
in general \emph{not} proportional to the identity. Why does this matter for
choosing $D$ in the moment expansion?
\end{exercise}

\begin{exercise}
Let $N=3$, $X$ nilpotent with $X^2\neq0$, $X^3=0$. Construct $p,\nu$ with
$M_0=M_1=0$, $M_2\neq0$, and verify Theorem~\ref{thm:krylov} by computing
$K_\Gamma$ explicitly from the adjugate.
\end{exercise}

\begin{exercise}
For the two-branch sector graph with one coupling insertion $V$, show
$M_0=0$ and $K_\Gamma=O(\Gamma^{-2})$, and derive the two visibility branches
$C_{\mathrm{abs}}=O(\Gamma^{-3})$ ($|\beta S_2/\gamma_g|\ll1$) and
$O(\Gamma^{-2})$ ($\gg1$).
\end{exercise}

\begin{exercise}[NV thermal collapse]
Assume phonon-induced orbital hopping with $\Gamma(T)\propto T^5$ at low
temperature. Using the Class III scaling of the NV spin-$\Lambda$, predict
how the EIT contrast scales with $T$ in each visibility branch, and propose a
plot (contrast vs $x(T)=\Gamma(T)/\Delta_{\mathrm{sel}}$) that would falsify
the theory.
\end{exercise}

\begin{exercise}[Symmetry audit]
Write down the audit table of Theorem~\ref{thm:symm} for the NV zero-field
spin-$\Lambda$ and identify which physical terms (spin--orbit? hyperfine?
strain?) break the spin sector. Conclude why the honest classification is
Class III rather than Class II.
\end{exercise}

% ==================================================================
\section*{Take-home messages}
% ==================================================================

\begin{keyidea}[Summary]
\begin{enumerate}
\item Classify \textbf{channels (configurations)}, not materials.
\item An optical dark vector ($\ker\Omega$), a stationary Lindblad dark
      state, a coherent pathway ($K_{12}K_{21}$), and observable EIT are
      \textbf{four different things}.
\item The correct no-go object is the sector-resolved correction
      $\dchi\equiv0$; a zero of the \emph{total} susceptibility is
      \textbf{perfect EIT}, not a no-go.
\item Fast full-rank damping generates a \textbf{moment hierarchy}: the first
      nonzero moment fixes the suppression power of $1/\Gamma$; finitely many
      vanishing moments certify an all-orders exact zero
      (Krylov/Cayley--Hamilton).
\item Rb: go. NV zero-field spin-$\Lambda$: asymptotic dissipative no-go,
      lifted quadratically by $B_\perp$. SiV/SnV orbital-$\Lambda$: kernel go
      whose observability is a budget question ($T$, strain, splitting).
\end{enumerate}
\end{keyidea}

\end{document}
